\section{\tr{General}}

\subsection{\tr{Preliminary Remarks}}

\subsubsection{\tr{Important Notices}}

\tr{Before using any part of the system, please read and understand this manual. All information in this
document is subject to change without notice.}


\subsubsection{\tr{Conventions and Pictograph Definitions}}

\tr{The safety instructions in SteFly operating manuals are the result of risk evaluations and hazard
analyses. In this document, the following hazard levels and information are considered:}

\picnote
	{pictograph-yellow-warning.pdf}
	{\tr{Pay special attention to critical notes marked with a yellow caution symbol, because
	non-observance may result in damage or any other critical situation.}}

\picnote
	{pictograph-red-caution.pdf}
	{\tr{A red caution symbol signalises that non-observance may result in injuries.}}

\picnote
	{pictograph-blue-cloud.pdf}
	{\tr{A blue cloud indicates useful information or tips.}}

\subsection{Safety}

\subsubsection{Safety \tr{Precautions}}
\picnote
	{pictograph-yellow-warning.pdf}
	{\tr{Duty to inform: Each person involved in the installation or operation of LARUS must read and observe the 
	safety-related parts of these operating instructions.}}


\subsubsection{\tr{Proper Use}}
\tr{LARUS Vario Display visualises data which are measured and calculated by the LARUS sensor box. LARUS
was designed to calculate direction and strength of thermals and wind quickly and reliably.
Therefore, the sensor unit combines data from high-precision sensors and GNSS receivers in
sophisticated algorithms.

LARUS Vario Display shall be installed in the instrument panel.

LARUS Gliding Sensor Unit is an additional feature to supply glider pilots with accurate information
about wind, vertical air movement as well as additional attitude of the aircraft. Its use is limited
to day VFR conditions. Security decisions must be made regardless of having installed LARUS.}


\subsubsection{\tr{Improper Use}}

\tr{Improper use will cause all claims for liability and guarantees to be forfeited. Improper use is
deemed to be all use for purposes deviating from those mentioned above, especially:}
\begin{itemize}
	\item 	\tr{Using LARUS data in non-VFR conditions or during night is forbidden. LARUS is not
			certified. Although LARUS provides AHRS data to XCSoar you must not rely on the
			artificial horizon display.}
	\item 	\tr{Using LARUS data during aerobatics or during flight conditions with high angle of attack
			(stall) or high g-forces. The algorithm was optimised for normal flight conditions.}
	\item 	tr{Operating LARUS Vario Display outside the operation conditions defined in technical data
			section, e.g. input voltage, temperature and humidity.}
\end{itemize}


\subsection{\tr{Quick Start Manual}}

LARUS Vario Display may be operated out of the box. Simply perform the following actions:
\begin{enumerate}
	\item 	\tr{Disassemble both rotary knobs with a 1,5 mm HEX key (scope of delivery)}
	\item 	\tr{Mount LARUS Vario Display in a 57 mm cutout of the instrument panel with 3 screws. It
			may be mounted in all four directions.}
	\item 	\tr{Assemble both rotary knobs}
	\item 	\tr{Connect the CAN port of LARUS Vario Display and LARUS Sensor Box with a 1:1 patch cable
			included in scope of delivery.}
	\item 	\tr{Power on LARUS.}
	\item 	\tr{Check that the satellite symbol on the screen is yellow or green and the current heading
			is shown.}
	\item 	\tr{Choose or create a proper polar for your glider.}
	\item 	\tr{Your LARUS Vario Display is ready to fly!}	
\end{enumerate}
